\documentclass[11pt]{article}
\usepackage[a4paper,margin=26mm]{geometry}
\usepackage{amsmath,amssymb,amsthm,booktabs,array}
\usepackage[colorlinks=true,linkcolor=blue,citecolor=blue,urlcolor=blue]{hyperref}
\setlength{\emergencystretch}{3em}
\newtheorem{theorem}{Theorem}
\newtheorem{lemma}[theorem]{Lemma}
\newtheorem{corollary}[theorem]{Corollary}
\newcommand{\norm}[1]{\left\lVert#1\right\rVert}
\newcommand{\upstream}[2]{\href{https://github.com/the-omega-institute/trureturing/blob/66d225eda69544653f1ca0e66cf3aea530b1a2c6/D5/S3/Combinatorics/Apollonian/#1.lean}{\texttt{#2}}}
\title{MathOverflow 516405, revision 1:\\
the second adjacency eigenvalue of Apollonian networks\\
and the fixed-rank extreme spectrum}
\author{Haobo Ma\\\small ChronoAI Pte Ltd, Singapore
\and Wenlin Zhang\\\small National University of Singapore\\\small ChronoAI Pte Ltd, Singapore}
\date{November 2, 2026}
\begin{document}
\maketitle
\begin{abstract}
We give an affirmative answer to the second-adjacency-eigenvalue question in
Rad Sarar Anando's MathOverflow question 516405, revision 1:
$\theta_2(G_m)/2^{m/2}\to\sqrt2$ for deterministic Apollonian networks
formed by subdividing every inner face and leaving the outer face unchanged.
The same argument determines every fixed-rank eigenvalue at both ends of the
ordinary adjacency spectrum. The squared limiting positive values are $3$,
three copies of $2$, and $3^j$ copies of $3\cdot2^{-j}$ for each $j\ge1$.
An exact cutoff Gram identity gives a finite-stage error bound and an
$O_r(2^{-m/4})$ convergence rate at each fixed rank $r$.
The original second-eigenvalue limit has a pinned Lean declaration; the
fixed-rank extension and error estimates are proved here in writing.
\end{abstract}
\begin{center}\small
\copyright\ 2026 Haobo Ma and Wenlin Zhang.
Licensed under \href{https://creativecommons.org/licenses/by/4.0/}{CC BY 4.0}.
\end{center}

\section{The question and the extreme eigenvalues}
MathOverflow question 516405, revision 1, posted by Rad Sarar Anando on
October 26, 2026 \cite{anando}, asks whether
$\theta_2(G_m)/2^{m/2}\to\sqrt2$. Here $G_0$ is a plane $K_4$ with outer
vertices $a,b,c$ and centre $o$. To obtain $G_{m+1}$, insert one vertex
in each inner triangular face of $G_m$ and join it to the face's three
corners. Retain every old edge and never subdivide the outer face.
Write $A_m$ for ordinary unweighted adjacency, $N_m=|V(G_m)|$, and
\[
\theta_1(G_m)\ge\theta_2(G_m)\ge\cdots\ge\theta_{N_m}(G_m)
\]
for its real eigenvalues with multiplicity. All matrix norms below are
Euclidean operator norms.

Let $d_1\ge d_2\ge\cdots$ enumerate the multiset
\begin{equation}\label{eq:weights}
\{3\}\;\cup\;\{2,2,2\}\;\cup\;
\bigcup_{j\ge1}\{3\cdot2^{-j}\text{ repeated }3^j\text{ times}\}.
\end{equation}
Its ordering is unambiguous because the values after $2$ strictly decrease.

\begin{theorem}\label{thm:limits}
For every fixed positive integer $r$,
\begin{equation}\label{eq:limits}
\frac{\theta_r(G_m)}{2^{m/2}}\longrightarrow\sqrt{d_r},
\qquad
\frac{\theta_{N_m+1-r}(G_m)}{2^{m/2}}\longrightarrow-\sqrt{d_r}.
\end{equation}
Both convergences have error $O_r(2^{-m/4})$.
In particular, the first four positive limits are
$\sqrt3,\sqrt2,\sqrt2,\sqrt2$, and ranks $5,6,7$ have limit $\sqrt{3/2}$.
\end{theorem}

This also answers the source's question about other leading eigenvalues.
Frieze and Tsourakakis \cite[Theorem 2]{frieze} prove the analogous
square-root-of-degree relation with high probability for random Apollonian
networks, where one face is chosen at each insertion. Their random model
is different from simultaneous subdivision of all inner faces.
Earlier spectral work includes Andrade and Miranda \cite{andrade}, and
Cardoso, Andrade and Souza \cite{cardoso}, whose uniform tight-binding
operator is proportional to adjacency. Serva's exactly solvable model
\cite{serva} instead regularizes the outer-face asymmetry and uses
degree-normalized hopping. Its spectrum does not give the constants in
\eqref{eq:limits}. The original second-eigenvalue limit is covered by
the pinned Lean declaration \cite{source}. The complete cutoff matrix
argument below gives the fixed-rank and quantitative consequences.

\section{Counts and a uniform remainder bound}
Initially there are three inner faces, and every subdivision replaces one
inner face by three. Thus there are $3^{m+1}$ inner faces at stage $m$,
exactly $3^j$ vertices are born at stage $j\ge1$, and
\begin{equation}\label{eq:degrees}
N_m=\frac{3^{m+1}+5}{2},\qquad
\deg_m(o)=3\cdot2^m,\qquad
\deg_m(a)=\deg_m(b)=\deg_m(c)=2^{m+1}+1.
\end{equation}
Indeed, every inner face incident with a persistent vertex gives two
incident child faces and adds one neighbour. The centre starts with three
incident inner faces, and each outer corner starts with two. A vertex
born at stage $j\ge1$ starts with degree and inner-face count both equal
to three, so its degree at stage $m\ge j$ is
\begin{equation}\label{eq:innerdegree}
\deg_m(v)=3\cdot2^{m-j}.
\end{equation}
These recurrences prove all the displayed counts directly.

Fix $0\le k<m$, put $H=V(G_k)$ and $L=V(G_m)\setminus H$, and order
the adjacency matrix by $H,L$. No later operation adds edges between
persistent vertices. Hence
\begin{equation}\label{eq:split}
A_m=P_{m,k}+\operatorname{diag}(A_k,C_{m,k}),\qquad
P_{m,k}=\begin{pmatrix}0&B_{m,k}\\B_{m,k}^{T}&0\end{pmatrix},
\end{equation}
where $B_{m,k}$ is old-to-late adjacency and $C_{m,k}$ is adjacency
induced on $L$.

\begin{lemma}\label{lem:norms}
For $0\le k<m$,
\begin{equation}\label{eq:norms}
\norm{A_k}\le6\,2^{k/2},\qquad
\frac{\norm{C_{m,k}}}{2^{m/2}}
\le3\sqrt2\,2^{-k/2}.
\end{equation}
\end{lemma}
\begin{proof}
If an undirected graph admits an orientation with outdegrees at most
$p$ and indegrees at most $q$, its adjacency matrix $C$ satisfies
\[
|x^TCx|\le2\sum_{u\to v}|x_u x_v|
\le2\Big(\sum_{u\to v}x_u^2\Big)^{1/2}
       \Big(\sum_{u\to v}x_v^2\Big)^{1/2}
\le2\sqrt{pq}\,\norm{x}^2.
\]
Since $C$ is symmetric, $\norm{C}\le2\sqrt{pq}$.
Orient every nonseed edge from its newer to its older endpoint.
Every nonseed vertex has precisely three older neighbours, and there
are no edges within a nonseed birth generation.
In $G_k$, order the four seed vertices and orient their edges toward
the earlier vertex. All outdegrees are then at most three. By
\eqref{eq:degrees}--\eqref{eq:innerdegree}, all indegrees are at most
$3\cdot2^k$, including $k=0$. This gives the first bound.

On $L$, all births are at least $k+1$, so every degree, and hence every
indegree, is at most $3\cdot2^{m-k-1}$. The outdegree bound remains
three. Therefore
$\norm{C_{m,k}}\le6\,2^{(m-k-1)/2}$; dividing by $2^{m/2}$ gives
the second bound.
\end{proof}

\section{An exact cutoff Gram identity}
For $v\in H$ let $f_k(v)$ be its incident inner-face count at stage $k$.
Let $F_k$ be the symmetric matrix with zero diagonal whose entry
$F_k(u,v)$, for $u\ne v$, is the number of inner faces containing both
$u$ and $v$. Put
\[
D_k=\operatorname{diag}\bigl(2^{-k}f_k(v):v\in H\bigr),\qquad
Q_{m,k}=2^{-m}B_{m,k}B_{m,k}^T.
\]
The entries of $D_k$ are $3$ for $o$, $2$ for each outer vertex, and
$3\cdot2^{-j}$ for a vertex born at stage $j\le k$.

\begin{lemma}\label{lem:gram}
Writing $t=m-k$, one has
\begin{equation}\label{eq:gram}
Q_{m,k}=(1-2^{-t})D_k+t\,2^{-m}F_k,
\qquad
\norm{Q_{m,k}-D_k}\le3(1+2t)2^{-t}.
\end{equation}
\end{lemma}
\begin{proof}
The diagonal entry of $B_{m,k}B_{m,k}^T$ at $v$ counts its neighbours
born after $k$. The successive increments are
$f_k(v),2f_k(v),\ldots,2^{t-1}f_k(v)$, summing to
$(2^t-1)f_k(v)$.

For distinct $u,v\in H$, its $(u,v)$ entry counts their common late
neighbours. A new common neighbour is born precisely in a face
containing both vertices. Such a face is a triangle, so $u,v$ must
already be adjacent. Under subdivision, exactly one child of that face
still contains both $u$ and $v$. The number of faces containing the
pair is therefore unchanged at each later stage. Each of the $t$
subdivisions contributes exactly $F_k(u,v)$ common neighbours, giving
the off-diagonal entry $tF_k(u,v)$. These identities prove the matrix
formula in \eqref{eq:gram}.

Every inner face incident with $v$ has exactly two other corners.
Consequently the row sum of $F_k$ at $v$ is $2f_k(v)$.
The degree recurrences give $f_k(v)\le3\cdot2^k$, so the symmetric
nonnegative matrix $F_k$ has $\norm{F_k}\le6\cdot2^k$.
For completeness, a symmetric nonnegative matrix with row sums at most
$M$ satisfies
$|x^TF_kx|\le\sum_{u,v}F_k(u,v)(x_u^2+x_v^2)/2\le M\norm{x}^2$.
Finally $\norm{D_k}=3$, whence
\[
\norm{Q_{m,k}-D_k}\le3\,2^{-t}+6t\,2^{k-m}
=3(1+2t)2^{-t}.
\]
\end{proof}

We use two elementary spectral facts. For symmetric matrices of the same
size, decreasing eigenvalues satisfy
\begin{equation}\label{eq:weyl}
|\lambda_r(X+E)-\lambda_r(X)|\le\norm{E}.
\end{equation}
Indeed,
$\lambda_r(X)=\max_{\dim U=r}\min_{x\in U,\,\norm{x}=1}x^TXx$:
the span of the first $r$ eigenvectors gives one inequality, and every
$r$-dimensional subspace intersects the span of eigenvectors numbered
$r,\ldots,n$ nontrivially, giving the other.
On the unit sphere $|x^TEx|\le\norm{E}$, which proves
\eqref{eq:weyl}.

Second, the nonzero eigenvalues of
$\left(\begin{smallmatrix}0&B\\B^T&0\end{smallmatrix}\right)$ are
$\pm\sqrt\mu$, where $\mu$ ranges over the positive eigenvalues of
$BB^T$, with multiplicity. If $BB^Tv=\mu v$, the vectors
$(v,\pm B^Tv/\sqrt\mu)$ give these eigenvalues, and the eigenvector
equations give the converse. Symmetry makes geometric and algebraic
multiplicities equal. If $B$ has $q$ rows, $\ell$ columns and rank $p$,
the remaining $q+\ell-2p$ eigenvalues are zero. Since $p\le\ell$, the
first $q$ decreasing eigenvalues of this dilation are the square roots
of the $q$ decreasing eigenvalues of $BB^T$, including zeros.
The last $q$ eigenvalues, read upward from the bottom, are their
negatives.

\section{Finite-stage errors and the limiting spectrum}
The first $N_k$ entries of the infinite sequence \eqref{eq:weights}
are exactly the decreasing diagonal entries of $D_k$.

\begin{theorem}\label{thm:error}
For $0\le k<m$, $1\le r\le N_k$, and $t=m-k$, define
\begin{equation}\label{eq:error}
E_{m,k,r}=\max\{6\,2^{-t/2},\,3\sqrt2\,2^{-k/2}\}
+\frac{3(1+2t)2^{-t}}{\sqrt{d_r}}.
\end{equation}
Then
\[
\left|\frac{\theta_r(G_m)}{2^{m/2}}-\sqrt{d_r}\right|\le E_{m,k,r},
\qquad
\left|\frac{\theta_{N_m+1-r}(G_m)}{2^{m/2}}+\sqrt{d_r}\right|
\le E_{m,k,r}.
\]
\end{theorem}
\begin{proof}
Let $\mu_r$ be the $r$th decreasing eigenvalue of the positive
semidefinite matrix $Q_{m,k}$. Lemma~\ref{lem:gram} and
\eqref{eq:weyl} give
$|\mu_r-d_r|\le\delta_t:=3(1+2t)2^{-t}$.
Since $\mu_r\ge0$ and $d_r>0$,
\[
|\sqrt{\mu_r}-\sqrt{d_r}|
=\frac{|\mu_r-d_r|}{\sqrt{\mu_r}+\sqrt{d_r}}
\le\frac{\delta_t}{\sqrt{d_r}}.
\]
The dilation identity identifies the first and last $r$th eigenvalues
of $2^{-m/2}P_{m,k}$ with $\sqrt{\mu_r}$ and $-\sqrt{\mu_r}$.
The perturbation in \eqref{eq:split} is block diagonal, with norm
exactly $\max\{\norm{A_k},\norm{C_{m,k}}\}$.
After division by $2^{m/2}$, Lemma~\ref{lem:norms} bounds it by the
maximum term in \eqref{eq:error}. Apply \eqref{eq:weyl} at each
end of the spectrum and use the triangle inequality.
\end{proof}

\begin{proof}[Proof of Theorem~\ref{thm:limits}]
Choose $k$ with $N_k\ge r$ and let $m\to\infty$ in
Theorem~\ref{thm:error}. Each upper limit of the two nonnegative
errors is at most $3\sqrt2\,2^{-k/2}$. Letting $k\to\infty$
proves \eqref{eq:limits} along all natural stages.
For the rate, instead take $k=\lfloor m/2\rfloor$, so $N_k\ge r$
for every sufficiently large $m$. The maximum term is
$O(2^{-m/4})$, and the final term is
$O_r(m2^{-m/2})=O_r(2^{-m/4})$. This gives the asserted rate.
\end{proof}

\begin{corollary}\label{cor:degrees}
Let $\Delta_r(G_m)$ be the $r$th largest degree, with multiplicity.
For every fixed $r$,
\[
\frac{\theta_r(G_m)}{\sqrt{\Delta_r(G_m)}}\longrightarrow1,
\qquad
\frac{\theta_{N_m+1-r}(G_m)}{\sqrt{\Delta_r(G_m)}}\longrightarrow-1.
\]
\end{corollary}
\begin{proof}
For $m\ge2$ the degrees are ordered by the centre, the three outer
vertices, and then increasing birth generation. Their formulas in
\eqref{eq:degrees}--\eqref{eq:innerdegree} give
$2^{-m}\Delta_r(G_m)\to d_r>0$ for every fixed $r$.
Divide the two limits in Theorem~\ref{thm:limits} by
$\sqrt{2^{-m}\Delta_r(G_m)}$.
\end{proof}

\section{Source coverage and verification}
The following coverage refers to the deterministic ordinary adjacency
operator and the unchanged outer face throughout.
\begin{center}\small
\begin{tabular}{@{}p{0.47\linewidth}p{0.46\linewidth}@{}}
\toprule
Item in MathOverflow 516405, revision 1 & Coverage here\\
\midrule
Question: second eigenvalue limit $\sqrt2$ & Theorem~\ref{thm:limits}, $r=2$.\\
Equivalent upper-limit formulation & Full limit proved, hence the upper limit equals $\sqrt2$.\\
Observation: $\theta_1^2$ follows the central degree & Corollary~\ref{cor:degrees}, $r=1$.\\
Question about other leading eigenvalues & Every fixed rank, both spectral ends; Theorems~\ref{thm:limits} and~\ref{thm:error}.\\
Random Apollonian square-root pattern & Prior probabilistic theorem \cite[Theorem 2]{frieze}; a different graph model.\\
Posted finite-stage bounds and numerical values & Not premises of the proof; the bounds in Theorem~\ref{thm:error} are derived directly.\\
\bottomrule
\end{tabular}
\end{center}

The declaration
\upstream{SecondEigenvalueLimit}{SecondEigenvalueLimit.result}
formalizes precisely the original second-largest ordinary adjacency
limit for the inner-face graph construction, along every natural stage
and with eigenvalue multiplicity.
The companion files have the following roles:
\upstream{InnerFaceStages}{InnerFaceStages} defines the finite graph,
births and faces and derives its degree and common-neighbour laws;
\upstream{CutoffSpectralBounds}{CutoffSpectralBounds.late\_adjacency\_normalized\_le}
gives the normalized late-block estimate;
\upstream{CutoffGramLimit}{CutoffGramLimit.normalized\_gram\_tendsto}
gives fixed-cutoff Gram convergence; and
\upstream{CutoffOrderedSpectrum}{CutoffOrderedSpectrum.cutoff\_second\_eq\_sqrt\_gram}
gives the second Gram-to-dilation identity under its positive-Gram-eigenvalue
hypothesis. The fixed-rank limits, exact identity \eqref{eq:gram},
finite-stage estimate \eqref{eq:error}, rate and degree corollary have
written proofs above; they are not asserted as additional Lean declarations.

The accompanying \texttt{verify.py} checks graph counts, all vertex degrees,
face incidence, the exact integer cutoff Gram identity, and directed
degree bounds for stages $0$ through $8$ and $26$ cutoff--stage pairs
with cutoffs at most $3$. The largest graph has $9{,}844$ vertices.
It also computes up to eight eigenvalues at each spectral end; the
maximum absolute eigenpair residual is below $2.2\cdot10^{-11}$.
These finite checks test the formulas; the infinite limits follow
from the written proof. The argument concerns fixed ranks, not ranks
growing with $m$, and does not identify a full empirical spectral
distribution or a normalized-adjacency or Laplacian spectrum.

\begin{thebibliography}{9}
\bibitem{anando}
Rad Sarar Anando,
\emph{Second adjacency eigenvalue of Apollonian networks: is
$\theta_2(G_m)/2^{m/2}\to\sqrt2$?}
MathOverflow question 516405, revision 1, October 26, 2026.
\href{https://mathoverflow.net/questions/516405}{Question};
\href{https://mathoverflow.net/revisions/8C9D6CD5-E3A3-4B25-BDB6-FF5A9A3D3A2B/view-source}{revision source}.

\bibitem{andrade}
R. F. S. Andrade and J. G. V. Miranda,
Spectral properties of the Apollonian network,
\emph{Physica A} \textbf{356} (2005), 1--5.
\href{https://doi.org/10.1016/j.physa.2005.05.003}{doi:10.1016/j.physa.2005.05.003}.

\bibitem{cardoso}
A. L. Cardoso, R. F. S. Andrade and A. M. C. Souza,
Localization properties of a tight-binding electronic model on the
Apollonian network, \emph{Phys. Rev. B} \textbf{78} (2008), 214202.
\href{https://doi.org/10.1103/PhysRevB.78.214202}{doi:10.1103/PhysRevB.78.214202};
\href{https://arxiv.org/abs/0809.4087v1}{arXiv:0809.4087v1}.

\bibitem{frieze}
Alan Frieze and Charalampos E. Tsourakakis,
Some properties of random Apollonian networks,
\emph{Internet Mathematics} \textbf{10} (2014), 162--187.
\href{https://doi.org/10.1080/15427951.2013.796300}{doi:10.1080/15427951.2013.796300}.
Preprint: \emph{High degree vertices, eigenvalues and diameter of random
Apollonian networks},
\href{https://arxiv.org/abs/1104.5259v3}{arXiv:1104.5259v3}, Theorem 2.

\bibitem{serva}
Maurizio Serva,
Exactly solvable tight-binding model on the RAN: fractal energy spectrum
and Bose--Einstein condensation,
\emph{Journal of Statistical Mechanics: Theory and Experiment}
(2014), P08018.
\href{https://doi.org/10.1088/1742-5468/2014/08/P08018}{doi:10.1088/1742-5468/2014/08/P08018};
\href{https://arxiv.org/abs/1402.4661v1}{arXiv:1402.4661v1}.

\bibitem{source}
The trureturing project,
\emph{Second eigenvalue of the inner-face Apollonian stages}.
\href{https://github.com/the-omega-institute/trureturing/blob/66d225eda69544653f1ca0e66cf3aea530b1a2c6/D5/S3/Combinatorics/Apollonian/SecondEigenvalueLimit.lean}{\texttt{SecondEigenvalueLimit.result}},
\texttt{SecondEigenvalueLimit.lean}, revision
\texttt{66d225eda69544653f1ca0e66cf3aea530b1a2c6}.
\end{thebibliography}
\end{document}
